\documentclass[11pt]{article}

\usepackage[margin=1in]{geometry}
\usepackage{graphicx} % Required for inserting images
\graphicspath{{../manybodyscars/figures/}}
\usepackage{hyperref}
\usepackage{mathtools,amssymb}
\usepackage{amsthm}
\usepackage{mathrsfs}
\usepackage{fancyhdr}
\usepackage{subcaption}
\captionsetup{font=small,labelfont=bf}
\usepackage{float}
\usepackage{mathdots}
\usepackage{physics}
\usepackage{indentfirst}
% \usepackage{fontspec} % Enable with XeLaTeX or LuaLaTeX for custom fonts.

\hypersetup{
  colorlinks=true,
  linkcolor=blue,
  citecolor=blue,
  urlcolor=blue
}

\pagestyle{fancy}
\fancyhf{}
\fancyhead[L]{From Quantum Many-Body Scars to Slow Dynamics}
\fancyhead[R]{Research note}
\fancyfoot[C]{\thepage}
\setlength{\headheight}{14pt}

\newcommand{\Hpxp}{H_{\mathrm{PXP}}}
\newcommand{\Ztwo}{\mathbb{Z}_2}
\newcommand{\Xtilde}{\widetilde{X}}

\title{From Quantum Many-Body Scars to Slow Dynamics\\
  \large Exact Diagonalization of a Deformed PXP Chain}
\author{Zihao Xu}
\date{\today}

\begin{document}

\maketitle

\begin{abstract}
This note summarizes a numerical study of a periodically closed, constrained spin
chain derived from the PXP model. The Hamiltonian combines the standard
Rydberg-blockaded spin-flip term with staggered one-spin and three-spin diagonal
perturbations. Exact diagonalization is used to resolve the spectrum, identify
eigenstates with anomalously large overlap with the charge-density-wave product
state $\ket{\Ztwo}$, and compute equilibrium thermodynamic quantities. The
largest-overlap scarred eigenstate is found in the low-energy part of the finite-size
spectrum. The nonequilibrium calculation probes a connected local
$X$-autocorrelation after a thermal quench. Its early-time decay is fitted directly,
without selecting peaks, by linear regression of $\log C_X(t)$ against $\log t$.
The observed power-law decay provides a concrete measure of slow relaxation. It
motivates comparison with Griffiths-like dynamics, but it does not by itself
establish a Griffiths phase.
\end{abstract}

\section{Motivation}

Experiments on Rydberg-atom arrays revealed persistent oscillations following a
quench from an ordered state, despite the absence of conventional integrability.
The corresponding constrained spin model supports atypical
eigenstates embedded in an otherwise thermal spectrum. These states, now called
quantum many-body scars, violate the usual expectation that all finite-energy-density
eigenstates of a generic interacting system locally resemble thermal equilibrium.

The present project asks how this scarred spectral structure changes when the PXP
chain is deformed so that the most strongly scarred state moves toward low energy.
The numerical workflow connects three complementary diagnostics: the exact energy
spectrum and $\Ztwo$ overlaps, the canonical heat capacity, and the relaxation of a
local autocorrelation function. The first diagnoses atypical eigenstates, the second
tests their thermodynamic context, and the third measures the timescale on which local
memory is lost.

\section{Constrained spin model}

Consider a periodic chain of $L$ spin-$1/2$ degrees of freedom. Let $X_i$ and
$Z_i$ denote Pauli operators, and define
\begin{equation}
  n_i=\frac{1+Z_i}{2},
  \qquad
  P_i=1-n_i=\frac{1-Z_i}{2}.
\end{equation}
The Rydberg blockade restricts the Hilbert space to configurations satisfying
\begin{equation}
  n_i n_{i+1}=0
  \quad \text{for every }i,
\end{equation}
with site labels understood modulo $L$. Within this constrained space, a local spin
flip is equivalent to
\begin{equation}
  \Xtilde_i=P_{i-1}X_iP_{i+1},
\end{equation}
and the standard PXP Hamiltonian is
\begin{equation}
  \Hpxp=\sum_{i=0}^{L-1}\Xtilde_i.
\end{equation}

The model studied here adds staggered diagonal perturbations,
\begin{equation}
  H(g,r)=\Hpxp
  +g\sum_{i=0}^{L-1}(-1)^i
  \left(Z_i+r Z_{i-1}Z_iZ_{i+1}\right).
  \label{eq:hamiltonian}
\end{equation}
The parameter $g$ sets the overall strength of the deformation and $r$ fixes the
relative weight of the three-spin term. The staggering reduces one-site translation
symmetry to translation by two sites. Equation~\eqref{eq:hamiltonian} is implemented
in a custom constrained basis and diagonalized with QuSpin.

The reference product state is
\begin{equation}
  \ket{\Ztwo}=\ket{1010\cdots},
\end{equation}
which is invariant under translation by two sites and therefore belongs to the
zero-momentum sector of that symmetry.

\section{Exact diagonalization and scarred eigenstates}

For the spectral calculation, the Hamiltonian is represented in the two-site
translation sector $k=0$. The current calculation uses $L=20$ and obtains the full
dense eigendecomposition
\begin{equation}
  H(g,r)\ket{E_n}=E_n\ket{E_n}.
\end{equation}
For each eigenstate, the principal scar diagnostic is its overlap with the ordered
state,
\begin{equation}
  w_n=\abs{\braket{\Ztwo}{E_n}}^2.
  \label{eq:overlap}
\end{equation}
The spectral location of the initial state can also be summarized by
\begin{equation}
  E_{\Ztwo}=\mel{\Ztwo}{H}{\Ztwo}
  =\sum_n E_n w_n.
\end{equation}

The numerical spectrum contains a small set of eigenstates whose values of $w_n$ are
anomalously large compared with nearby eigenstates. These outliers are identified as
quantum many-body scar candidates. For the deformation of principal interest,
$g=-0.4$ and $r=0.2$, the eigenstate with the largest $\Ztwo$ overlap lies in the
low-energy region of the finite-size spectrum. This placement differs from the
familiar picture of scarred states near the middle of the spectrum and makes
low-temperature thermodynamic diagnostics especially relevant.

The half-chain entanglement entropy of every eigenstate is evaluated after lifting the
states from the symmetry-reduced basis to the full constrained basis. It is used as
an additional color-coded spectral diagnostic. Unfolded nearest-neighbor level
spacings are computed from the middle two thirds of the spectrum, where edge effects
are weaker. These quantities help distinguish isolated nonthermal states from a
global loss of ergodicity, although a systematic finite-size scaling analysis is
required before drawing a thermodynamic-limit conclusion.

Figure~\ref{fig:spectrum} collects the spectral diagnostics and the corresponding
quench dynamics. The upper-left panel makes the anomalously large $\Ztwo$ overlaps
visible and places the strongest-overlap states toward the low-energy side of the
spectrum. The lower panels show that the same initial state produces long-lived,
damped oscillations in both a local observable and the return probability.

\begin{figure}[H]
  \centering
  \includegraphics[width=0.96\linewidth]
  {pxp_spectrum_L=20_g=-0.4_r=0.2.png}
  \caption{Spectral and dynamical diagnostics for $L=20$, $g=-0.4$, and
  $r=0.2$. Upper left: overlap $\abs{\braket{\Ztwo}{E_n}}^2$ versus energy;
  the color scale gives the half-chain entanglement entropy, and the dashed line
  marks $E_{\Ztwo}$. Upper right: distribution of unfolded nearest-neighbor level
  spacings. Lower panels: the site-averaged nearest-neighbor correlation
  $L^{-1}\sum_i\expval{Z_iZ_{i+1}}$ and the return probability
  $\abs{\braket{\Ztwo}{\psi(t)}}^2$ after a quench from $\ket{\Ztwo}$.}
  \label{fig:spectrum}
\end{figure}

\section{Canonical energy and heat capacity}

The thermodynamic calculation uses the eigenvalues from all two-site momentum
sectors, rather than only the $k=0$ sector used for the scar plot. If
$E_{n,k}$ denotes the spectrum in an independent momentum block, the partition
function is reconstructed as
\begin{equation}
  Z(\beta)=\sum_{k\in\mathcal{K}_{\mathrm{ind}}}m_k
  \sum_n e^{-\beta E_{n,k}},
  \label{eq:partition}
\end{equation}
where $m_k=2$ for a conjugate pair $k$ and $-k$, while self-conjugate sectors have
$m_k=1$. In practice, the global ground-state energy is subtracted inside the
Boltzmann factors to avoid numerical underflow; this shift is restored in the mean
energy.

With $k_{\mathrm B}=1$, the canonical internal energy and heat capacity are
\begin{align}
  U(T) &= \expval{H}_\beta
  =\frac{1}{Z}\sum_{k,n}m_k E_{n,k}e^{-\beta E_{n,k}},\\
  C_V(T) &= \frac{\partial U}{\partial T}
  =\beta^2\left(\expval{H^2}_\beta-\expval{H}_\beta^2\right).
  \label{eq:capacity}
\end{align}
The temperature scan in Figure~\ref{fig:capacity-T} uses $L=20$, $r=0.2$,
$0.02\leq T\leq2$, and $g=0,-0.1,\ldots,-0.5$. The internal energy rises
smoothly with temperature. Each heat-capacity curve has a broad maximum, whose
position shifts to higher temperature as the deformation becomes stronger. These
maxima show where thermal fluctuations sample the largest concentration of spectral
weight.

\begin{figure}[H]
  \centering
  \includegraphics[width=\linewidth]{pxp_capacity_T_L=20_r=0.2.png}
  \caption{Canonical thermodynamics as a function of temperature for $L=20$ and
  $r=0.2$. The left panel shows the total internal energy $U=\expval{H}$; the
  right panel shows the total heat capacity $C_V$. Each curve corresponds to a
  different deformation strength $g$.}
  \label{fig:capacity-T}
\end{figure}

Figure~\ref{fig:capacity-g} presents the complementary scan in which temperature is
fixed and $g$ varies from $-0.5$ to $0$. At the lowest temperatures, the internal
energy follows the deformation of the low-energy spectrum while $C_V$ remains close
to zero. At higher temperatures, the heat capacity becomes sensitive to $g$ across
the full interval. Taken together, Figures~\ref{fig:capacity-T} and
\ref{fig:capacity-g} show how the same spectral rearrangement appears along two cuts
through the $(g,T)$ plane. A broad maximum or shoulder in $C_V$ is a crossover
indicator at these finite sizes, not by itself evidence of a phase transition.

\begin{figure}[H]
  \centering
  \includegraphics[width=\linewidth]{pxp_capacity_g_L=20_r=0.2.png}
  \caption{Canonical thermodynamics as a function of deformation strength for
  $L=20$ and $r=0.2$. The left and right panels show $U(g)$ and $C_V(g)$,
  respectively, at fixed temperatures $T=0.01$, $0.1$, $0.2$, $0.5$, $0.8$,
  and $1$.}
  \label{fig:capacity-g}
\end{figure}

\section{Connected local autocorrelation}

The dynamical calculation starts from a canonical state of the undeformed PXP
Hamiltonian,
\begin{equation}
  \rho_0=\frac{e^{-\beta\Hpxp}}
  {\Tr(e^{-\beta\Hpxp})},
\end{equation}
and evolves operators with the deformed Hamiltonian $H_1=H(g,r)$,
\begin{equation}
  \Xtilde_i(t)=e^{iH_1t}\Xtilde_i e^{-iH_1t}.
\end{equation}
The measured quantity is the site-averaged connected autocorrelation
\begin{equation}
  C_X(t)=\frac{1}{L}\sum_{i=0}^{L-1}\operatorname{Re}\left[
  \Tr\!\left(\rho_0\Xtilde_i(t)\Xtilde_i\right)
  -\Tr\!\left(\rho_0\Xtilde_i(t)\right)
   \Tr\!\left(\rho_0\Xtilde_i\right)
  \right].
  \label{eq:correlation}
\end{equation}
The subtraction removes the disconnected one-point contribution. Because $\rho_0$
need not commute with $H_1$, the time-dependent one-point function is retained in
that subtraction. The code evaluates Eq.~\eqref{eq:correlation} in the eigenbasis of
$H_1$, allowing all requested times to be obtained from precomputed matrix
contractions and phase factors.

The current parameter set is $L=16$, $g=-0.4$, $r=0.2$, and $T=1$. Times are
sampled geometrically over $0.1\leq t\leq100$, which provides uniform resolution on
a logarithmic time axis.

\subsection{Power-law fit}

The observed decay is modeled by
\begin{equation}
  C_X(t)=A t^{-\alpha}.
  \label{eq:powerlaw}
\end{equation}
For positive correlation values, Eq.~\eqref{eq:powerlaw} becomes the linear relation
\begin{equation}
  \log C_X(t)=\log A-\alpha\log t.
  \label{eq:logfit}
\end{equation}
Ordinary least-squares regression is applied directly to all correlation samples in
the early-time interval
\begin{equation}
  0.1\leq t\leq t_{\mathrm{fit}},
  \qquad t_{\mathrm{fit}}=50.
\end{equation}
No peak detection or upper-envelope selection enters the fit. The fitted exponent is
minus the regression slope, and the amplitude is the exponential of the intercept.
The plot uses logarithmic axes and marks the fitted time interval with a translucent
shaded region without an edge line.

Figure~\ref{fig:autocorrelation} shows the correlation data and the fitted power law.
The regression gives $\alpha=0.486$ for this parameter set. Oscillations remain around
the fitted trend, but the dashed line captures the overall decrease across the shaded
time window. Data beyond $t_{\mathrm{fit}}$ are displayed for comparison and do not
enter the regression.

\begin{figure}[H]
  \centering
  \includegraphics[width=0.88\linewidth]
  {pxp_autocorr_avg_L=16_g=-0.4_r=0.2_T=1.0.png}
  \caption{Connected local $X$-autocorrelation for $L=16$, $g=-0.4$,
  $r=0.2$, and $T=1$ on logarithmic axes. Orange points inside the shaded region
  are the data used in the regression over $0.1\leq t\leq50$. The green dashed
  line is the fitted power law $C_X(t)=At^{-\alpha}$ with $\alpha=0.486$; data at
  later times are excluded from the fit.}
  \label{fig:autocorrelation}
\end{figure}

The approximately linear behavior in the log--log representation shows that the
finite-size autocorrelation decays as a power law over the selected window. Such slow
relaxation is consistent with the broad phenomenology often associated with
Griffiths regimes near localization transitions. The present model,
however, is clean and staggered rather than spatially random. Power-law decay over a
finite time window is therefore a motivation for further tests, not a sufficient
identification of a Griffiths phase. Useful next steps include varying $L$, the fit
window, temperature, and deformation parameters; comparing alternative observables;
and testing whether the fitted exponent approaches a stable large-system value.

\section{Summary}

This study follows a single chain of evidence from spectral structure to equilibrium
and nonequilibrium observables. Exact diagonalization reveals scarred eigenstates
with anomalously large $\Ztwo$ overlap and places the strongest-overlap state in the
low-energy part of the finite-size spectrum. A symmetry-resolved reconstruction of
the full partition function yields the internal energy and heat capacity. Finally,
the connected local $X$-autocorrelation displays a power-law decay when the early-time
data are fitted directly on logarithmic axes. Together, these calculations provide
a controlled finite-size framework for studying how scar physics may coexist with,
or evolve into, slow many-body relaxation.

\end{document}
